\label{implementation:strict}

The delivered \code{code/global_rewrite.py} is a complete standalone
checker for the seven-rule calculus. Its input is an exact four-key JSON
object containing a nonnegative integer \code{variables}, a list
\code{context} of retained formulas, a formula \code{goal}, and a
primitive \code{proof} tree. It validates the full tagged Eq/Q/D language
and numerical variable scopes. The input contract includes open sequents
with the declared numerical parameters, as well as closed theorems.
Acceptance requires a normal exit and exactly \code{ACCEPT} followed by
one newline. The program imports only Python's standard JSON and system
modules. Neither a producer nor a theorem database is trusted.

Numerical terms are encoded as \code{['v',i]}, \code{['z]},
\code{['s',a]} and \code{['i',n,a,h]}; the final argument of the iterator
binds one variable. Formulas are \code{['=',s,t]}, \code{['q',B]} and
\code{['d',B]}. Variables use de Bruijn indices, with index zero the most
recent binder. The body of D has its case variable at index zero and
its witness variable at index one. The RW instruction contains only its
two premise proofs: \code{['rw',p,q]}. It contains no numerical or formula
context to be instantiated.

\subsection{One traversal with three precise modes}

The helper \code{walk(a,t,s,k,d)} uses the same structural recursion for
lifting, substitution and global replacement. Here $d$ records the number
of binders traversed. When a source pattern is supplied, a node matching
the source lifted through $d$ binders is replaced by the target lifted
through those binders. The traversal stops at that match. An unmatched
variable below depth $d$ remains bound; an unmatched index at least $d$
is changed by $k$. The supported calls are:

\begin{center}\small
\begin{tabular}{@{}lll@{}}
\toprule
Operation & Source pattern & Index offset\\
\midrule
Lifting by $k$ & none & $k$\\
Substitution, dropping one variable & variable zero & $-1$\\
Successor instance under a retained eigenvariable & variable zero & $0$\\
Global replacement & the proved source term & $0$\\
\bottomrule
\end{tabular}
\end{center}

The matching test precedes descent into children, which enforces the stated
top-down replacement convention. In particular, newly inserted material
is not visited again. The pattern and image are ordinary checked terms;
the internal value \code{None} denotes the absence of a pattern and is
not numerical syntax.

For DI, the matrix is first instantiated at the supplied witness. The
progress hypothesis is retained exactly when the matrix is an equation
and the witness is the literal zero term; otherwise the retained
hypothesis is its zero case. In each branch, the ambient context is lifted
when the eigenvariable is introduced. DE similarly lifts the context and
the claimed conclusion before checking its continuation. This implements
freshness through the scope discipline rather than by a convention imposed
on the certificate producer.

\subsection{Whole-program measurements}

The original X comparator, \code{reference/one_context_seven.py}, is
\emph{our matched seven-rule adaptation} of the incoming
\code{one_context/direct} checker: its two introduction constructors are
replaced by the already established witness-sensitive W introduction.
It is not represented as an unchanged published 3,181-byte file. The
callback and unified-traversal controls retain that X rule and change only
the traversal. Relative to the unified X control, the new program changes
only the equality-rule branch.

All rows charge an entire standalone module, including syntax validation,
scope checks, inference checking, JSON input and the verdict driver.
Actual docstrings are removed, and CPython 3.12.14
\code{ast.unparse} plus one final newline gives the normalized source.
Significant-token counts omit comments, encoding and whitespace markers;
AST counts include every node of \code{ast.walk}. Gzip uses level 9,
\code{mtime=0}, and zlib 1.3.2, including the gzip wrapper.

\begin{center}\small
\begin{tabular}{@{}p{0.46\linewidth}rrrr@{}}
\toprule
Complete standalone module & Bytes & Tokens & AST & Gzip-9\\
\midrule
One-context X: original traversal & 3181 & 1029 & 1089 & 976\\
One-context X: callback traversal & 3215 & 1059 & 1124 & 983\\
One-context X: unified traversal & 3206 & 1044 & 1096 & 987\\
Template-free global replacement & 3147 & 1027 & 1072 & 983\\
\bottomrule
\end{tabular}

\end{center}

The new module is 34 normalized bytes, two significant tokens and
17 AST nodes smaller than the original-traversal X control, a byte
reduction of approximately $1.07\%$. Its gzip size is seven bytes larger;
compression is reported as supplementary information. With exactly the
same unified traversal on both sides, removing the supplied context and
its checking work saves 59 bytes, 17 tokens and 24 AST nodes.

These are modest matched implementation improvements. The compact
arity-code frontier uses a different certificate language and input
contract. The table does not establish a new corpus-wide minimum, and
does not provide a lower bound on possible checkers. The mathematical
contribution is the elimination of the context parameter, together with
the large-marker completeness proof; code size is measured separately.

\subsection{Executed checks and their limits}

The command \code{python3 tests/verify_strict.py} regenerates all
certificates and the measurement table. The authoritative receipts are
\code{evidence/strict_validation.json} and
\code{evidence/strict_metrics.json}; the latter records source hashes,
raw and normalized byte counts, and exact runtime versions.

There are 38 named fixtures: 14 valid derivations and 24 expected
rejections. Four standalone programs, each run under ordinary Python and
\code{python -O}, give 304 process checks, with zero mismatches. The X
controls receive explicit X translations of the R proofs. Their retained
contexts, variable counts and claimed conclusions are identical. A
separate fixed-seed audit compares 600 liftings, 200 ordinary
substitutions and 200 substitutions without dropping the eigenvariable
against the original implementation, and checks 200 root replacements:
1,200 further comparisons, all matching.

Positive cases include large-marker congruence with accidental passive
occurrences of the source, simultaneous replacement in both equation
endpoints, nested iterator binders, a closed congruence theorem, and
equation induction that uses its hypothesis. The independent logical
naming construction exercises DI, DE and QE as well. Rejection controls
cover bound-variable capture, the wrong top-down convention, a case
hypothesis substituted for a required progress hypothesis, incorrect
computation endpoints, wrong premise kinds, malformed syntax, Boolean
and negative indices, and escaping witnesses.

During review, a preliminary callback traversal failed to preserve an
unmatched variable leaf: it attempted to descend into the variable's
integer payload. The nested-binder fixtures detected the defect. That
prototype and its preliminary size are not the delivered result. The
unified traversal handles every variable explicitly, and the final
receipts refer to the corrected source hashes.

The implementation audit also checks the inference invariants directly:
each checked premise has a well-scoped formula; lifting preserves the
bound/free distinction; substitution removes exactly its designated
variable when requested; and replacement inserts a term only after
lifting it through the encountered binders. These are written source
arguments supported by executable cases. The finite tests are not a
proof-assistant verification of the Python implementation or of
bi-interpretability. As for the baseline, the universal proof-system
statement abstracts from finite process memory and recursion limits;
a resource failure is never acceptance.
